% ============================================================
% Formulario, argomento 18: legge dei grandi numeri e teorema centrale del limite (dispensa, capitolo 18)
% ============================================================
\argomento{Legge dei grandi numeri e teorema centrale del limite}{Dispensa: capitolo 18}[Grandi numeri e limite
  centrale]

\voce{iid}{Variabili i.i.d.}{indipendenti e identicamente distribuite: il modello delle prove ripetute}
  {f_{X_1, \ldots, X_n}(x_1, \ldots, x_n) = \prod_{i=1}^{n} f_X(x_i)}
  {\variante{discrete}{p_{X_1, \ldots, X_n}(x_1, \ldots, x_n) = \prod_{i=1}^{n} p_X(x_i)}
   \variante{stessa media e stessa varianza}{E[X_i] = \mu, \quad \mathrm{VAR}(X_i) = \sigma^2 \quad \text{per ogni } i}}

\voce{mediacampione}{Media campione}{come la media campionaria \vv{mediacamp}; lettere piccole: i dati osservati}
  {\overline{x}_n = \frac{1}{n} \sum_{i=1}^{n} x_i \qquad \overline{X}_n = \frac{1}{n} \sum_{i=1}^{n} X_i}
  {\variante{dati i.i.d.}{E\left[ \overline{X}_n \right] = \mu, \quad \mathrm{VAR}\left[ \overline{X}_n \right] =
     \frac{\sigma^2}{n}}
   \variante{deviazione standard}{\sigma_{\overline{X}_n} = \frac{\sigma}{\sqrt{n}}}}

\voce{lgn}{Legge debole dei grandi numeri}{dati i.i.d. con media $\mu$ e varianza $\sigma^2$}
  {\lim_{n \to \infty} P\left( \abs{\overline{X}_n - \mu} > \epsilon \right) = 0 \qquad \forall\, \epsilon > 0}
  {\variante{con Chebyshev, per ogni $n$}{P\left( \abs{\overline{X}_n - \mu} > \epsilon \right) \leq
     \frac{\sigma^2}{n\, \epsilon^2}}
   \variante{per avere probabilità al più $\delta$ basta}{n \geq \frac{\sigma^2}{\delta\, \epsilon^2}}}

\voce{frequenzeprob}{Le frequenze tendono alle probabilità}{$f_n(E)$: frazione delle $n$ prove in cui si verifica $E$}
  {\lim_{n \to \infty} P\big( \abs{f_n(E) - P(E)} > \epsilon \big) = 0 \qquad \forall\, \epsilon > 0}
  {\variante{con l'indicatore di $E$ nella prova $i$}{X_i \in \{0, 1\}, \quad E[X_i] = P(E), \quad f_n(E) =
     \overline{X}_n}
   \variante{$N_i$ volte il simbolo $a_i$ in $n$ dati, $p_i = p_X(a_i)$}{P(N_i = k) = \binom{n}{k}\, p_i^k
     (1 - p_i)^{n - k}}
   \variante{frequenza relativa}{E\left[ \frac{N_i}{n} \right] = p_i, \quad \mathrm{VAR}\left[ \frac{N_i}{n} \right]
     = \frac{p_i (1 - p_i)}{n}}}

\voce{convergenze}{Modi di convergere}{$Y_n$ successione di variabili aleatorie, $c$ un numero}
  {\text{in probabilità:} \quad \lim_{n \to \infty} P\left( \abs{Y_n - c} > \epsilon \right) = 0 \quad \forall\,
   \epsilon > 0}
  {\variante{quasi certa (con probabilità uno)}{P\left( \lim_{n \to \infty} Y_n = c \right) = 1}
   \variante{in media quadratica}{\lim_{n \to \infty} E\left[ (Y_n - c)^2 \right] = 0}
   \variante{le ultime due implicano la prima}{}
   \variante{dalla media quadratica, con Markov}{P\left( \abs{Y_n - c} > \epsilon \right) \leq \frac{E\left[ (Y_n -
     c)^2 \right]}{\epsilon^2}}
   \variante{legge forte dei grandi numeri, dati i.i.d.}{P\left( \lim_{n \to \infty} \overline{X}_n = \mu \right) = 1}}

\voce{sommaiid}{Somma di $n$ variabili i.i.d.}{ciascuna con media $\mu_X$ e varianza $\sigma_X^2$}
  {Z_n = \sum_{i=1}^{n} X_i: \qquad \mu_{Z_n} = n\, \mu_X \qquad \sigma_{Z_n}^2 = n\, \sigma_X^2}
  {\variante{somma standardizzata, con media 0 e varianza 1}{S_n = \frac{Z_n - n\, \mu_X}{\sqrt{n}\, \sigma_X}}}

\voce{tcl}{Teorema centrale del limite}{$X_i$ i.i.d. con varianza finita, qualunque sia la loro legge}
  {\lim_{n \to \infty} P(S_n \leq x) = \Phi(x) = \int_{-\infty}^{x} \frac{1}{\sqrt{2\pi}}\, e^{-\frac{t^2}{2}} \dd t
   = 1 - Q(x)}
  {\variante{per $n$ grande, la somma}{Z_n \simeq \mathcal{N}\left( n\, \mu_X,\ n\, \sigma_X^2 \right)}
   \variante{e la media campione}{\overline{X}_n \simeq \mathcal{N}\left( \mu,\ \frac{\sigma^2}{n} \right), \quad
     P\left( \abs{\overline{X}_n - \mu} > \epsilon \right) \simeq 2\, Q\left( \frac{\epsilon \sqrt{n}}{\sigma}
     \right)}
   \variante{binomiale con $n$ grande}{\mathcal{B}(n, p) \simeq \mathcal{N}\big( np,\ np(1 - p) \big)}
   \variante{convergenza rapida al centro, lenta sulle code; non vale per la Cauchy}{}}

\voce{generatrice}{Funzione generatrice dei momenti}{lo strumento della dimostrazione del teorema \oltre}
  {M_Y(t) = E\left[ e^{tY} \right]}
  {\variante{gaussiana standard}{M(t) = e^{\frac{t^2}{2}}}
   \variante{somma di variabili indipendenti}{M_{Y_1 + \cdots + Y_n}(t) = M_{Y_1}(t) \cdots M_{Y_n}(t)}
   \variante{per $t$ piccolo}{M_Y(t) = 1 + t\, E[Y] + \frac{t^2}{2}\, E[Y^2] + \cdots}
   \variante{la stessa $M$ vicino a $t = 0$ dà la stessa legge; non esiste per tutte le variabili (la Cauchy)}{}}
